\begin{proof}[Proof of Proposition~\ref{shear:prop:disjoint-transfer}]
Write \(A=\widehat f(\rho)\) and abbreviate \(S_i(D)\) to \(S_i\).
If \(z=0\), then \(f=0\) and every assertion is immediate. Assume \(z>0\).

\paragraph{The common small-type cover.}
Disjoint supports embed \(H_1\times\cdots\times H_r\) in \(G\):
the factors commute, and a product that is identity restricts to the
identity on every \(M_i\). Complete reducibility and the
direct-product description of irreducibles
\cite[Theorems~3.12 and~4.25 and Proposition~3.14]{etingof} give
\[
 V_\rho\big|_{H_1\times\cdots\times H_r}
 =
 \bigoplus_{\boldsymbol\tau}
 (V_{\tau_1}\otimes\cdots\otimes V_{\tau_r})
       \otimes\mathcal A_{\boldsymbol\tau}.
\]
The sum is over distinct occurring product types, with their full
multiplicity spaces. Each has
\(\prod_i\dim\tau_i\le D\), so some factor has degree at most
\(D^{1/r}\). Unused labels can change the occurring types and their
multiplicities, but do not change this product bound.

Let \(P_{i,\tau}\) be the orthogonal projection onto all copies of
the \(H_i\)-type \(\tau\), and put
\[
 P_i=\sum_{\substack{\tau\in\widehat H_i\\
                         \dim\tau\le D^{1/r}}}P_{i,\tau},
 \qquad Q_i=P_i\prod_{j<i}(I-P_j).
\]
On each product-type summand, \(P_i\) is identity or zero according
as the \(i\)th factor is small or large. Hence the \(P_i\) commute,
and at least one is identity on each summand. Therefore
\begin{equation}
 I\preceq\sum_{i=1}^rP_i,\qquad
 \sum_{i=1}^rQ_i=I,\qquad Q_iQ_j=0\quad(i\ne j).
 \label{shear:eq:disjoint-projector-cover}
\end{equation}
The \(Q_i\) assign a whole product type to its first small factor;
the \(P_i\) may overlap. Both families retain every multiplicity.
For every \(u>0\), the definition of \(C_u\) gives
\[
 S_i\le D^{(u+2)/r}
       \sum_{\tau\in\widehat H_i}(\dim\tau)^{-u}
 \le C_uD^{(u+2)/r}.
\]

\paragraph{The orbit estimate.}
For \(v\in V_\rho\), write \(v_i=Q_iv\) and
\(W_i=\Span\{\rho(h)v_i:h\in H_i\}\). The single-vector orbit lemma
\cite[Lemma~\Lref{l:single-orbit}]{partial-fourier} gives
\(\dim W_i\le S_i\). Indeed, in all copies of a type \(\tau\), the
orbit of its component is contained in the image of
\[
 \operatorname{End}(V_\tau)\longrightarrow
 V_\tau\otimes\mathcal A_\tau,\qquad
 T\longmapsto(T\otimes I)v_{i,\tau},
\]
whose dimension is at most \((\dim\tau)^2\), independently of the
multiplicity.

Let \(\Pi_i\) project onto \(W_i\). Since \(W_i\) is \(H_i\)-invariant,
the projection onto \(\rho(g)W_i\) depends only on \(gH_i\).
The coset cap and Schur averaging yield
\[
 \sum_g f(g)\|\Pi_{\rho(g)W_i}w\|^2
 \le K\frac{\dim W_i}{D}\|w\|^2.
\]
Here the uniform average of
\(\rho(g)\Pi_i\rho(g)^*\) is \((\dim W_i/D)I\), by its trace and
irreducibility. Weighted Cauchy--Schwarz retains the mass \(z\):
\[
 |\langle w,Av_i\rangle|^2
 \le z\|v_i\|^2\sum_g f(g)\|\Pi_{\rho(g)W_i}w\|^2
 \le\frac{zKS_i}{D}\|w\|^2\|v_i\|^2.
\]
The \(v_i\) are orthogonal. Summing and applying Cauchy--Schwarz over
\(i\) gives
\[
 |\langle w,Av\rangle|
 \le\left(\frac{zK}{D}\sum_iS_i\right)^{1/2}\|w\|\|v\|.
\]
This proves \eqref{shear:eq:noncovering-orbit-general}.
Its specialization \(u=2\), with \(z\le1\), is
\eqref{shear:eq:orbit-bound}; the condition \(r>4\) makes that
dimension exponent negative.

\paragraph{The isotypic estimate.}
For a subgroup \(H\le G\) and an irreducible \(H\)-type \(\tau\) of
degree \(a\), define
\[
 e_{H,\tau}(h)=\frac{a}{|H|}\overline{\chi_\tau(h)}
 \quad(h\in H),\qquad e_{H,\tau}=0\quad\text{off }H.
\]
Its transform is the projection \(P_{\sigma,H,\tau}\) onto all copies
of \(\tau\) in \(\sigma|_H\). Thus
\(\widehat{f*e_{H,\tau}}(\sigma)=
\widehat f(\sigma)P_{\sigma,H,\tau}\): a bound on the left cosets
\(gH\) controls a projection on the right of the Fourier matrix.
The companion's cosetwise isotypic lemma
\cite[Lemma~\Lref{l:central-isotypic}]{partial-fourier} applies to
every subgroup \(H\). If \(f(gH)\le K[G:H]^{-1}\), it gives
\begin{equation}
 \sum_{\sigma\in\widehat G}D_\sigma
       \|\widehat f(\sigma)P_{\sigma,H,\tau}\|_{\HS}^2
 =|G|\sum_g|(f*e_{H,\tau})(g)|^2\le Kza^2.
 \label{shear:eq:noncovering-central}
\end{equation}
This is the central-projection step; it retains the whole projected
Fourier matrix and the mass factor.

Trace the first inequality of \eqref{shear:eq:disjoint-projector-cover}
against \(A^*A\), and apply \eqref{shear:eq:noncovering-central} to
each selected \(H_i\)-type after discarding the other ambient
irreducibles. Then
\[
 D\|A\|_{\HS}^2
 \le D\sum_i\sum_{\dim\tau\le D^{1/r}}
       \tr(A^*AP_{i,\tau})
 =D\sum_i\sum_{\dim\tau\le D^{1/r}}
       \|AP_{i,\tau}\|_{\HS}^2
 \le Kz\sum_iS_i.
\]
The trace inequality uses positivity of
\(A(\sum_iP_i-I)A^*\); it does not require \(A^*A\) to commute with
the projections. This proves
\eqref{shear:eq:noncovering-isotypic-general} using the overlapping
\(P_i\), rather than the orbit assignment \(Q_i\).

\paragraph{The coefficient alternative.}
The independent coefficient calculation in
\cite[Lemma~\Lref{l:isotypic-alternatives}]{partial-fourier}
bounds each coset separately. Its norm and orientation give the
selected-type refinement in \eqref{shear:eq:hs-general}.
For one subgroup \(H\le G\), put \(F=|G|f\), so
\(\E_{h\sim U_H}F(xh)=[G:H]f(xH)\le K\).
For an \(H\)-type \(\tau\) of degree \(a\), choose one representative
\(x\) of each coset and set
\[
 T_x=\E_{h\sim U_H}F(xh)\tau(h),\qquad
 q(xh)=a\sum_{j,\ell=1}^a(T_x)_{j\ell}\overline{\tau(h)_{j\ell}}.
\]
The function \(q\) on this coset is the orthogonal projection of
\(F\) onto the conjugate coefficient space of \(\tau\). Schur
orthogonality and the nonnegative weights \(F(xh)\) give
\begin{equation}
 \E_{h\sim U_H}|q(xh)|^2
 =a\|T_x\|_{\HS}^2
 \le a^2\|T_x\|_{\op}^2
 \le a^2\bigl(\E_{h\sim U_H}F(xh)\bigr)^2
 \le K^2a^2.
 \label{shear:eq:coset-coefficient-norm}
\end{equation}
Expanding \(T_x\) and using, for \(h,k\in H\),
\(\overline{\chi_\tau(k^{-1}h)}=\chi_\tau(kh^{-1})\)
shows that these coset projections form the single function
\(q=|G|(f*e_{H,\tau})\). Thus their tested Fourier matrix is
\(\widehat f(\sigma)P_{\sigma,H,\tau}\), with the same right-hand
orientation as above. Averaging the preceding norm bound over cosets
and applying complex finite-group Plancherel gives
\begin{equation}
 \sum_{\sigma\in\widehat G}D_\sigma
       \|\widehat f(\sigma)P_{\sigma,H,\tau}\|_{\HS}^2
 =\E_{g\sim U_G}|q(g)|^2\le K^2a^2.
 \label{shear:eq:coefficient-one-type}
\end{equation}
The calculation counts a type once, regardless of its multiplicity,
and assumes no pointwise bound inside the coset. The same positive
trace comparison now gives
\(\|A\|_{\HS}^2\le K^2D^{-1}\sum_iS_i\).
Using \(S_i\le C_{10}D^{12/r}\) proves
\eqref{shear:eq:hs-general}, including its original sum over the
types of the unmodified subgroups \(H_i\).

\paragraph{The branching refinement.}
Each \(H_i\) is conjugate to the standard subgroup \(S_{|M_i|}\)
of \(S_N\). Young branching
\cite[Theorem~9.2]{James1978} shows that every occurring type
\(\tau\vdash |M_i|\) is a subdiagram of \(\lambda\); different
corner-removal sequences can give arbitrary multiplicity.
If \(\lambda_1\ge\lambda'_1\), the diagram \(\tau\) has at most
\(k_\lambda\) boxes below its first row. Its lower diagram, together
with its fixed size, determines the first row, so at most
\(M(k_\lambda)\) distinct types can occur. If
\(\lambda'_1>\lambda_1\), transpose both diagrams for this count.
Transposition preserves degrees and does not replace \(f\) or its
Fourier matrix by a sign-twisted law.

Nonoccurring isotypic projections are zero. In the positive trace
comparison, apply \eqref{shear:eq:noncovering-central} only to the
occurring selected types and use
\((\dim\tau)^2\le D^{2/r}\). There are at most \(M(k_\lambda)\) of
them for each support, giving \eqref{shear:eq:hs-branching}.
The same argument includes \(k_\lambda=0\), where the representation
is trivial or sign, \(D=1\), and \(M(0)=1\). It also includes all
unequal nonempty supports and every restriction multiplicity.
\end{proof}
